\section{Results}

The results are presented in two parts. We first validate the raw outputs of the BFM, establishing that its point predictions are accurate and that its two uncertainty estimates are well-characterised. We then demonstrate how predictions and their associated uncertainty propagate to Species Sensitivity Distributions (SSDs) and the derivation of hazardous concentrations, one of the practical application forms of observed and predicted ecotoxicity data in environmental protection, assessment and management.

\subsection{Prediction accuracy}
The evaluation of the BFM predictive accuracy considers the comparison of observed and predicted LC50 values. Each predicted value is the model's best estimate of the LC50 for an untested species--chemical combination, a data-driven form of read-across whose use is explored in detail by \citet{PosthumaImprovingHazard}. The BFM model achieves an RMSE of 0.70 log mg/L in grouped 5-fold cross-validation ($k=32$, $n_{\text{iter}}=2000$, $n_{\text{burn}}=100$), where each held-out fold contains chemical--species--duration triplets never seen during training. This represents an improvement on \citet{PosthumaImprovingHazard}, who reported an RMSE of 0.85, the improvement likely attributable to the incorporation of chemical and species features in the model. Direct comparisons to other works are more difficult to make, for reasons described in  \citet{schur2025comparability}. \citet{Viljanen2023} tested a number of different models which yielded RMSE's in the 0.88-1 range, albeit on a different dataset. Achieving the best possible performance in this prediction setting is not the primary concern of this study, and therefore we do not attempt to further compare RMSE performance. 

The prediction accuracy of the BFM is illustrated in Figure~\ref{fig:pred_vs_observed}, which plots model predictions for species--chemical pairs at 48-hour exposure against the average of the actual observations for those pairs. The strong correlation gives a visual impression of how well the model recovers measured toxicity. Some scatter is expected from the observations themselves: each plotted x-value may rest on a single test or on many (Figure~\ref{fig:rsd-distr}), so the less-tested points carry more measurement noise. Consistent with this, \citet{PosthumaImprovingHazard} found that the predicted--observed correlation improves substantially when restricted to more-tested, robust datapoints.

\begin{figure}[htbp]
    \centering
    \includegraphics[width=0.6\linewidth]{images/predicted_vs_measured_48h.png}
    \caption{Each point on the plot is the mean measured log(mg/L) toxicity (average where multiple measurements available) vs.\ the mean predicted log(mg/L) toxicity for a unique chemical--species pair at 48 hours.}
    \label{fig:pred_vs_observed}
\end{figure}

A further inspection of the residuals in Figure~\ref{fig:residuals} allows for a deeper understanding of the biases of the model. Panel~(a) shows a scatterplot of the residuals against the predicted toxicity value. Aside from a few extremely outlying points, the shape of the plot is relatively uniform and suggests calibrated model behaviour. Panel~(b) shows the residuals scattered against the measured observations, showing a negative diagonal slope, indicating regression to the mean: very high measured values tend to be under-predicted while very low values tend to be over-predicted.
Having established that the model predicts accurately on average, the key question for risk assessment is whether we can trust individual predictions, particularly for chemicals with little testing history for the chemical--species pair itself.
\begin{figure}[htbp]
    \centering
    \includegraphics[width=1\linewidth]{images/bias_analysis_48h.png}
    \caption{Residual diagnostic plots showing (a) residuals vs predicted values, (b) residuals vs measured values revealing regression to the mean, and (c)  residual distribution histogram.}
    \label{fig:residuals}
\end{figure}


\subsection{Validating the two uncertainty components}

Each model prediction carries two uncertainty components. Epistemic uncertainty, derived from the variance across the model's posterior samples, is interpreted as the model confidence is in its mean prediction and it shrinks as more data constrains the fit. Aleatoric uncertainty is a per-chemical variance term the model learns directly, capturing the irreducible spread between repeat tests: the biological and experimental variability intrinsic to the assay, which represents the residual noise floor of the measurement rather than a gap in the model's features. Together these give a total predictive variance for each prediction; across all predictions, aleatoric variance usually contributes more than epistemic (${\sim}67\%$ of the total), 
indicating the model is relatively confident in its mean predictions but predicts spread in individual measurements, consistent with the high inter-test variability documented in Figure~\ref{fig:rsd-distr}. We validate these estimates in three steps: first that the aleatoric term recovers the spread seen in genuine repeat measurements, second that the two components respond to data availability as a reducible and an irreducible quantity should, and third that the combined predictive intervals are well-calibrated.

\subsubsection{Aleatoric uncertainty matches the spread of repeat measurements}

Repeated tests of the same species--chemical pair vary substantially (Figure~\ref{fig:rsd-distr}), so a single test result is best understood as one draw from a distribution rather than a fixed ground truth. The aleatoric term is the model's estimate of the width of that distribution, and we first check that it matches the spread of genuine repeat measurements.

Table~\ref{tab:replicate_stats} assesses whether the model's per-chemical aleatoric
uncertainty matches the empirical spread of genuine repeat measurements. For each
chemical--species--duration triplet measured at least five times, we compare the
empirical standard deviation of its replicates with the model's predicted aleatoric
standard deviation, summarising agreement by the ratio of the two. Across every
replicate bin the predicted and empirical standard deviations are of the same magnitude
(pooled predicted-to-empirical ratios range from $1.21$ at 5--9 replicates to $0.95$ at
100+), indicating that the learned aleatoric term recovers the scale of irreducible
measurement noise rather than absorbing systematic model error.

At low replicate counts the predicted spread is somewhat wider than the empirical one: at 5–9 replicates the bin-mean predicted SD exceeds the empirical SD by a pooled factor of about 1.2, with a median per-triplet ratio of 1.5. This gap closes as replicates accumulate — the median ratio falls monotonically to 0.90 by the 100+ bin and the pooled ratio to 0.95 — and it has a straightforward explanation that does not imply miscalibration. The comparison's denominator, the empirical SD, is itself computed from only a handful of replicates, and for small samples the sample standard deviation is both noisy and biased downward as an estimator of the true σ. The ratio is therefore inflated at low replicate counts even when the predicted SD is unbiased, and it settles toward one as the empirical estimate stabilises. The wide interquartile range in the sparse bins ([0.90, 2.69] at 5–9 replicates, tightening to [0.89, 1.03] at 100+) is a direct signature of this sampling noise in the denominator rather than of scatter in the model's estimate. A smaller, second contribution is the shared prior on the per-chemical precisions, which applies mild shrinkage to sparsely-measured chemicals toward the population's typical noise level.


For the data-rich chemicals that dominate regulatory SSD construction, the predicted and
empirical spreads coincide closely (Pooled Ratio $1.12$ at 50--99 replicates and $0.95$
at 100+), so the uncertainty estimates are most reliable in precisely the regime where
they are used. From a regulatory standpoint, the mild over-prediction that remains for
data-poor chemicals is conservative: it widens the confidence intervals on derived
hazard concentrations rather than producing false precision.

% TODO (LP Comment 343): Leo asks "Is this correctly worded? Is this what you mean here to be over-estimated?" 

% TODO (LP Comment 346): Leo says: "I would reserve the word 'conservative' for the
% regulatory-use context. What is the mathematical formulation here, solely driven by
% insights in the data and the predictions?" 
\begin{table}[ht]
    \centering
    \caption{Calibration of the model's aleatoric uncertainty against the empirical
    spread of genuine repeat measurements. For every unique chemical--species--duration
    triplet measured at least five times, the empirical standard deviation of the
    replicates is compared with the model's predicted aleatoric standard deviation;
    triplets are grouped into bins by replicate count. \textbf{N~groups}: number of
    triplets in the bin.  \textbf{Median Ratio} and \textbf{Ratio IQR}: median and
    interquartile range of the per-triplet predicted-to-empirical SD ratio. \textbf{Mean
    Pred~SD} and \textbf{Mean Emp~SD}: bin-mean predicted and empirical standard
    deviations (log mg/L). A ratio above one indicates the model predicts a wider spread
    than is observed empirically.}
    \label{tab:replicate_stats}
\begin{tabular}{lccccc}
    \toprule
    \textbf{Replicates} & \textbf{N groups} & \textbf{Median Ratio} & \textbf{Ratio IQR} & \textbf{Mean Pred SD} & \textbf{Mean Emp SD} \\
    \midrule
    5--9   & 1,803 & 1.474 & [0.90, 2.69] & 0.470 & 0.387 \\
    10--19 & 658   & 1.449 & [0.93, 2.47] & 0.467 & 0.381 \\
    20--49 & 227   & 1.315 & [0.93, 2.00] & 0.482 & 0.400 \\
    50--99 & 30    & 1.091 & [0.91, 1.51] & 0.520 & 0.463 \\
    100+   & 6     & 0.897 & [0.89, 1.03] & 0.533 & 0.558 \\
    \bottomrule
\end{tabular}
\end{table}

% TODO (LP Comment 385): Leo asks whether this figure (uncertainty vs observations)
% should be renumbered and positioned much earlier possibly before Table 2 or
% directly after the predictive accuracy subsection. 

\subsubsection{Epistemic and aleatoric uncertainty respond differently to data}

The two uncertainty components behave differently as data accrues, which is examined at the chemical level by plotting out-of-fold epistemic and aleatoric standard deviations against the total number of observations per chemical (Figure~\ref{fig:uncert_vs_num_obs}). Epistemic uncertainty showed a negative association with observation count (Spearman $\rho = -0.308$), consistent with the expectation that additional data constrains the chemical's latent factors and reduces parametric uncertainty: the reducible component shrinks where more evidence is available. This is the uncertainty that any statistical or model carries but that most methods leave unquantified; the posterior ensemble makes it explicit by measuring how much the predictions of different plausible model fits disagree.

The aleatoric component, by contrast, shows essentially no association with observation count (Spearman $\rho = 0.087$). This is what one expects of a genuinely irreducible noise term: the intrinsic biological and experimental variability of a chemical's toxicity tests should not systematically diminish simply because the chemical has been tested more often. The parameter $\alpha_c$ is estimated per chemical and captures the residual variance around the model's mean — measurement noise together with the inter‑species and inter‑duration spread that the rank‑32 factorisation does not resolve, and any remaining chemical‑specific misfit. Its flatness in observation count follows from the shared prior on the per‑chemical precisions. Because that prior's rate is learned from the whole population of chemicals, its scale tracks the typical per‑chemical residual variance in the data: a chemical with few observations has its precision pooled toward this population‑typical noise level, whereas one with many observations has its precision determined by the scatter of its own residuals. The pooling target is the same whether a chemical is sparsely or heavily tested, a chemical's intrinsic variability has no reason to depend on how often it happens to have been tested, so the estimate inherits no systematic trend with observation count. It is this pooling toward a data‑derived centre, rather than toward a fixed predetermined scale, that keeps the relationship flat. Epistemic and aleatoric uncertainty therefore behave exactly as a reducible and an irreducible quantity should: the first shrinks as evidence accumulates, the second stays broadly constant.

\begin{figure}[htbp]
    \centering
    \includegraphics[width=1\linewidth]{images/uncertainty_vs_observations_48h.png}
    \caption{Relationship between data availability and predicted uncertainty. Each point represents one chemical at 48\,h exposure duration. The x-axis shows the total number of observations (across all durations) available for that chemical in the training set. Chemicals with fewer than three total observations were excluded to avoid cold-start artifacts in the cross-validated uncertainty estimates (N\,=\,660 excluded; N\,=\,1,722 retained). (a) Mean epistemic (reducible) standard deviation (SD), computed as the average across out-of-fold posterior predictive variance for each chemical's 48\,h observations. (b) Mean aleatoric (non-reducible) SD, derived from the per-chemical precision parameter $\alpha_c$. Both uncertainty components are expressed in log mg/L units.}
    \label{fig:uncert_vs_num_obs}
\end{figure}

Figure~\ref{fig:variance_composition} summarises this contrast by decomposing the total predictive variance across observation-count bins. Epistemic variance declines steadily as observations accumulate, whereas aleatoric variance stays broadly stable; its apparent increase in the highest-count bin reflects the particular chemicals that bin contains rather than a systematic trend with data. Because epistemic shrinks while aleatoric does not, the share of total variance attributable to aleatoric noise grows with the amount of data: epistemic uncertainty dominates for sparsely tested chemicals, and aleatoric uncertainty dominates for well-tested ones.

% =========================================================================
% PLACEHOLDER: new variance-composition figure (the stacked-bar + share plot
% from discussion). Decide later whether to (a) include it here alongside
% Figure~\ref{fig:uncert_vs_num_obs}, or (b) replace Figure~\ref{fig:uncert_vs_num_obs}
% with it. If replacing, merge the two epistemic/aleatoric paragraphs above
% with the summary paragraph here. If dropping, delete the paragraph above too.
% =========================================================================
\begin{figure}[htbp]
    \centering
    \includegraphics[width=\linewidth]{images/variance_composition_vs_nobs.png}
    \caption{Composition of total predictive variance by data availability. Chemicals are binned by their number of observations. (a)~Mean predictive variance per bin, split into aleatoric (irreducible noise) and epistemic (model uncertainty) components.}
    \label{fig:variance_composition}
\end{figure}

The aleatoric parameter should not be interpreted as irreducible measurement noise in isolation: because $\alpha_c$ is shared across all species and durations for a given chemical, it absorbs any residual variance the model's rank-32 factorisation does not capture. The total predictive uncertainty for a given (chemical, species, duration) triplet is the sum of the epistemic and aleatoric components.

\subsubsection{Combined predictive intervals are well-calibrated}

To assess whether the total predictive uncertainty is well-calibrated, we construct posterior predictive intervals at a range of nominal coverage levels and measure how often the true held-out observation falls within them. For each out-of-fold observation and each posterior sample~$s$, we draw from $\mathcal{N}(\hat{y}^{(s)}, 1/\alpha_c^{(s)})$, combining the epistemic variation across posterior samples with the aleatoric noise learned for that chemical. Figure~\ref{fig:calibration} plots the resulting empirical coverage against the nominal level. The model is slightly conservative at low-to-mid levels (a 50\% credible interval captures approximately 57\% of observations) and close to nominal in the 90--95\% range most relevant for regulatory applications (87.7\% and 91.5\% empirical coverage, respectively). The mean absolute calibration error across all levels is 0.043. The model thus errs slightly toward overcoverage at mid-range levels and shows only mild undercoverage in the high-coverage range that matters most for regulatory use, with empirical coverage falling two to four percentage points below nominal at the 90--95\% levels.
%CHECK VALUES AFTER MODEL UPDATE OR AFTER FINAL RE-TRAINING

\begin{figure}[htbp]
    \centering
    \includegraphics[width=0.55\linewidth]{images/calibration_curve.png}
    \caption{Posterior predictive calibration evaluated on 70{,}670 out-of-fold observations. For each nominal credible interval level, empirical coverage reports the fraction of observations whose true value falls within the interval. The dashed line indicates perfect calibration.}
    \label{fig:calibration}
\end{figure}

To give a visual idea of what these uncertainty estimates look like in practice, Figure~\ref{fig:chosen_predictions} shows model predictions for five randomly selected chemicals alongside the species for which they were tested in the dataset, allowing side-by-side comparison of test results with predictions. The dark blue bars represent the epistemic model uncertainty around the mean, and the light blue bars extend to the total uncertainty, with the additional extent beyond the dark blue bars representing the per-chemical aleatoric spread.

% TODO (LP Comment 350): Leo notes that for some chemicals (e.g. Alachlor, a
% herbicide with specific MoA), observations may be dominated by algae tests that
% fit well, while predictions for other taxonomic groups may fail. He suggests
% showing chemicals grouped by MoA (narcotic, herbicide, insecticide) rather than
% random selection, and discussing "split-SSD" by taxonomic group.

% TODO (LP Comment 332): Leo suggests using the terms "data deserts" and "data
% mountains" to describe sparse vs dense regions of the heatmap, as an accessible
% metaphor for a lay audience.

\begin{figure}[htbp]
    \centering
    \includegraphics[width=0.85\linewidth]{images/example_predictions_48h.png}
    \caption{Predicted vs observed toxicity with decomposed uncertainty for five randomly selected chemicals at 48\,h exposure. For each chemical, the data on the tested species are sorted vertically by predicted toxicity. Blue dots show the model's mean prediction; dark blue bars indicate $\pm 1.96$ epistemic SD (posterior parametric uncertainty), and light blue bars extend to $\pm 1.96$ total SD (adding the per-chemical aleatoric component). Red crosses show the mean observed toxicity, with red bars indicating $\pm 1.96$ SD of replicate measurements where available.}
    \label{fig:chosen_predictions}
\end{figure}


\subsection{Propagating uncertainty into Species Sensitivity Distributions}

Species Sensitivity Distributions (SSDs) are one of the main ways ecotoxicity 
data are put to use in environmental protection, assessment, and management. An 
SSD describes how sensitivity to a chemical is distributed across species, and regulatory quantities are read off it as hazardous concentrations 
(HCx): the HC5 of a NOEC-based SSD underpins protective concentration standards, 
and the HC20 of an EC10-based SSD is used in life-cycle impact assessment. We use 
LC50 here as a methodological prototype, but the SSD machinery is the same for any 
endpoint. Traditionally, an SSD is obtained by fitting a lognormal curve to the 
few observed values available for a chemical. Building on the model-predicted SSD 
approach of \citet{PosthumaImprovingHazard}, which replaces that curve fit with 
an empirical CDF of model predictions covering every species in the dataset, we 
construct SSDs directly from the BFM's predictions. Crucially, each prediction now 
carries a characterised uncertainty, and the Bayesian structure of the BFM lets us 
propagate that uncertainty through the SSD and into the hazardous concentrations 
derived from it, which a point-prediction SSD cannot do. We present the results in 
two views: the uncertainty on each individual species prediction, and the 
uncertainty on the SSD curve as a whole.
% TODO (LP Comment 393): Leo asks for more explanation of the "ensemble SSD" concept
% upfront, defined clearly so that a layman would understand it intuitively. 

% TODO (LP Comment 394): Leo suggests reordering: show the per-point uncertainty
% view (current Figure 8 / \ref{fig:uncertainty_simple}) FIRST, to explain the raw
% BFM output per individual species prediction, and THEN introduce the ensemble SSD
% (current Figure 7 / \ref{fig:ssd_comparison}) as the aggregated end-user view.


The per-species uncertainty is shown directly by plotting individual species predictions with their decomposed uncertainty bars (Figure~\ref{fig:uncertainty_simple}). For each species, the dark blue bars show $\pm 1.96$ aleatoric standard deviations, the constant per-chemical noise term, and the light blue bars extend to $\pm 1.96$ total standard deviations, the additional width reflecting the species-specific epistemic component. This per-point view reveals where in the SSD the uncertainty concentrates and highlights the distinct behaviour of the two uncertainty components: the aleatoric term is constant across species within a chemical but is somewhat larger for the sparsely tested Chlorfenprop-methyl (SD~$\approx 0.64$) than for Atrazine (SD~$\approx 0.40$), reflecting modestly greater residual variability for the less-studied compound. The epistemic term varies across species and is also larger for Chlorfenprop-methyl, differing by roughly a third between the two chemicals. It is important to note that these per-point error bars describe uncertainty around individual predictions, not around the SSD curve as a whole; the ensemble approach below is needed for the latter.

\begin{figure}[htbp]
    \centering
    \includegraphics[width=0.495\linewidth]{images/atrazine_novel_uncertainty.png}
    \includegraphics[width=0.495\linewidth]{images/chlorfenprop_methyl_novel_uncertainty.png}
    \caption{Model-predicted SSDs for Atrazine (left) and Chlorfenprop-methyl (right) with per-species uncertainty bars. Dark blue: $\pm 1.96$ aleatoric SD (the constant per-chemical noise term); light blue: $\pm 1.96$ total SD (epistemic + aleatoric). The aleatoric component is constant within each chemical but somewhat larger for the sparsely tested Chlorfenprop-methyl.}
    \label{fig:uncertainty_simple}
\end{figure}

% TODO (LP Comment 395): Leo asks whether these SSD scatter dots could have shapes
% and colours for different taxonomic groups.

Figure~\ref{fig:ssd_comparison} shows the ensemble of predicted SSDs for two contrasting chemicals, each overlaid on the traditional lognormal SSD fitted to the observed species. For Atrazine, a data-rich compound with 54 tested species and a specific Mode of Action (photosynthesis inhibitor, herbicide), the posterior ensemble forms a narrow band. For Chlorfenprop-methyl (also a herbicide), which has observations for only 7 species, the ensemble fans out considerably, reflecting the greater epistemic uncertainty associated with a sparsely tested chemical. This contrast demonstrates the key practical value of the approach: the width of the ensemble directly communicates how much the SSD, and any regulatory quantity derived from it, should be trusted. For Chlorfenprop-methyl this width is driven by epistemic uncertainty, which is reducible: additional testing of this chemical would narrow the band. For a data-rich chemical the residual width instead reflects irreducible measurement noise that further testing would not remove.

% TODO (LP Comment 392): Leo suggests expanding this to a four-panel figure:
% (A) data-rich narcotic chemical, (B) data-poor narcotic, (C) data-rich herbicide,
% (D) data-poor herbicide. Dots coloured by taxonomic group (green = plants,
% blue = arthropods, red = vertebrates) to show that colours are mixed across the
% SSD for narcotics but grouped/clustered for specific-MoA chemicals.

For each of the $S = 1{,}900$ post-burn-in Gibbs samples retained during training, we evaluate the full factorisation machine equation at $\theta^{(s)}$ to obtain a deterministic toxicity prediction $\hat{y}_i^{(s)}$ for every species~$i$ for each chemical of interest. Each set of predictions is sorted to produce a complete SSD curve, yielding an ensemble of $S$ internally consistent SSDs. Because each curve is derived from a single coherent parameter draw, the correlated epistemic structure across species is preserved: the entire SSD shifts and rescales as the shared parameters vary, rather than individual points scattering independently. Credible intervals for any functional metric derived from an SSD, such as the HC5 or HC20 (Hazardous Concentration for 5\% or 20\% of the species, respectively), can then be obtained directly from the empirical distribution of that quantity across the posterior ensemble.

\begin{figure}[htbp]
     \centering
     \begin{subfigure}[b]{0.48\textwidth}
         \centering
         \includegraphics[width=\textwidth]{images/ssd_uncertainty_atrazine_48h.png}
         \caption{Atrazine (54 tested species)}
         \label{fig:ssd_atrazine}
     \end{subfigure}
     \hfill
     \begin{subfigure}[b]{0.48\textwidth}
         \centering
         \includegraphics[width=\textwidth]{images/ssd_uncertainty_chlorfenprop_methyl_48h.png}
         \caption{Chlorfenprop-methyl (7 tested species)}
         \label{fig:ssd_chlorfenprop}
     \end{subfigure}
     \caption{Ensemble SSDs for Atrazine (left) and Chlorfenprop-methyl (right) at 48\,h exposure, overlaid on the traditional lognormal SSD. The shaded band is the BFM 95\% credible region and the solid navy line its ensemble median, both constructed from 1{,}900 posterior samples; the red curve is the traditional lognormal SSD fitted to the observed species (black points). The ensemble is narrow for the data-rich Atrazine, where it closely tracks the traditional fit, and wide for the sparsely tested Chlorfenprop-methyl, where the traditional fit still falls within the credible band.}
     \label{fig:ssd_comparison}
\end{figure}

Overlaying the traditional lognormal SSD on the same axes shows where the two methods agree and how firmly that agreement is pinned down. For Atrazine the ensemble is narrow and the traditional fit sits almost on top of the ensemble median, the two hazard estimates agreeing closely (BFM HC20 $= -0.33$ against a traditional value of $-0.41$ log mg/L) and the traditional value falling well within the model's 95\% credible interval, $[-0.50, -0.17]$. For Chlorfenprop-methyl the ensemble instead spans more than a log unit: the traditional value ($-2.79$) still falls within the model's credible interval, $[-2.99, -1.55]$, but that interval is wide enough that any hazard concentration inferred from only seven tested species is left largely unconstrained. In both cases the traditional and model-predicted curves are read in the same way; what the ensemble adds is an explicit statement of how well the underlying data determine the SSD, and with it the hazardous concentration derived from it.


\subsection{Uncertainty-aware hazardous concentrations}

SSD models are used for two main purposes,  to derive Hazardous Concentrations that are used as protective standards (which is based on selecting a regulatory maximum-accepted Y-value, and deriving a concentration threshold as resulting X-value), and to derive the Potentially Affected Fraction of species (Y-value) given an ambient pollutant concentration (X-value). In this section, we illustrate how the ensemble SSD can be used for the former, to estimate an HCx value from the resulting SSD.

% TODO (LP Comment 409): Leo asks to expand the HC20 forest plot (Figure 10) with
% Y-axis information stating MoA and n_tests for each chemical (e.g.,
% "gamma-Cyhalothrin; Insecticide, n=xxx"). This would show (a) that data-rich
% chemicals have closer agreement between BFM and traditional, and (b) how MoA
% influences the pattern.

The ensemble of posterior SSDs enables direct derivation of hazardous concentrations with credible intervals. Figure~\ref{fig:hc20_30_chems} shows BFM-derived HC20 values with 95\% credible intervals alongside HC20 values obtained from traditional lognormal curve-fitting for a selection of chemicals. Two patterns are apparent: (i)~many of the traditionally derived HC20 values fall outside the BFM's 95\% credible interval, and (ii)~the BFM systematically predicts lower (that is, more protective) HC20 values.

This conservative bias is confirmed across the full set of 412 chemicals for which both approaches could be computed (Figure~\ref{fig:hc20correlation}; Pearson $r = 0.949$). The BFM-derived HC20 values fall consistently below the 1:1 line, particularly for chemicals with fewer observed species (the darker points in Figure~\ref{fig:hc20correlation}); the shift below the line grows as the number of tested species falls (Spearman $\rho = 0.29$ between species count and the signed HC20 difference). The relationship is a statistical tendency rather than a rule: individual chemicals fall on both sides of the 1:1 line, but the balance is strongly toward more protective BFM estimates. The likely explanation is that traditional SSDs are fitted to the small, often non-representative subset of species that have been tested in the laboratory, whereas the BFM predictions cover all 1{,}267 species in the dataset, drawing on cross-species and cross-chemical read-across. The broader species coverage may include taxa that are more sensitive but have simply never been tested against a given chemical. From a regulatory perspective, this conservatism means the BFM approach is less likely to underestimate hazard, a desirable property, since an SSD fitted to only a handful of species carries large sampling uncertainty and, by omitting untested and potentially more sensitive taxa, can overestimate the protective concentration.

Taken together, the SSD and HC results show that the BFM provides not only a hazard estimate for each chemical but also an explicit statement of how much evidence supports it, which is the information a risk assessor needs to decide how far to rely on a model-derived value.

\begin{figure}[htbp]
    \centering
    \includegraphics[width=0.8\linewidth]{images/hc20_forest_plot.png}
    \caption{HC20 values derived from BFM posterior ensemble (blue, with 95\% credible intervals) compared to HC20 values from traditional lognormal curve-fitting (red markers). Chemicals are ordered by BFM-derived HC20.}
    \label{fig:hc20_30_chems}
\end{figure}

\begin{figure}[htbp]
    \centering
    \includegraphics[width=0.65\linewidth]{images/hc20_correlation_trad_vs_mc.png}
    \caption{Correlation between traditional (lognormal-fitted) and BFM-derived HC20 values for 412 chemicals with at least 5 tested species (Pearson $r = 0.949$). Points are coloured by the number of observed species per chemical. The systematic shift below the 1:1 line indicates more protective BFM estimates, particularly for data-poor chemicals.}
    \label{fig:hc20correlation}
\end{figure}
